\paragraph{二面体特化机制：由预解根到单二次域闭合}
在结论 \ref{con:xi-terminal-zm-s4-resolvent-cubic-discriminant-equals-delta} 与结论
\ref{con:xi-terminal-zm-resolvent-cubic-weierstrass-birational} 的同一记号体系下，
平面三次预解曲线 \(\mathscr R(z,y)=0\) 的 \(\QQ\)-有理点提供了一个可计算的特化机制：当 \(z\) 在基域中出现时，四次谱覆盖的伽罗瓦群从泛型的 \(S_4\) 下降到配对稳定子，且 Ferrari 分解所需的两处平方根可在同一二次扩张中闭合。

\begin{theorem}[函数域上的二面体特化：预解根强制 \texorpdfstring{$D_4$}{D4} 伽罗瓦群]\label{thm:xi-terminal-zm-dihedral-d4-specialization-functionfield}
令 \(K:=\QQ(E_{\mathscr R})\) 为椭圆曲线 \(E_{\mathscr R}\) 的函数域（结论 \ref{con:xi-terminal-zm-resolvent-cubic-weierstrass-birational}），并在 \(K\) 中取
\[
y=-\frac{8X}{Y+19+14X-X^2},
\qquad
z=1-\frac{2X(7-X)}{Y+19+14X-X^2}.
\]
则 \((y,z)\) 满足 \(\mathscr R(z,y)=0\)。
将四次谱多项式
\[
\Pi(\lambda,y)=\lambda^{4}-\lambda^{3}-(2y+1)\lambda^{2}+\lambda+y(y+1)
\in K[\lambda]
\]
视为 \(K\) 上的可分四次多项式，则其伽罗瓦群满足
\[
\mathrm{Gal}\bigl(\Pi(\lambda,y)/K\bigr)\ \subseteq\ D_4.
\]
此外，存在 \(P\in E_{\mathscr R}(\QQ)\) 使得特化多项式 \(\Pi(\lambda,y(P))\in\QQ[\lambda]\) 的伽罗瓦群同构于 \(D_4\)；
从而
\[
\mathrm{Gal}\bigl(\Pi(\lambda,y)/K\bigr)\ \cong\ D_4.
\]
因此存在无穷多个 \(P\in E_{\mathscr R}(\QQ)\) 使得 \(\Pi(\lambda,y(P))\) 的伽罗瓦群同构于 \(D_4\)。
\end{theorem}

\begin{proof}[证明要点]
由 \(\mathscr R(z,y)=0\) 可知，\(\Pi(\lambda,y)\) 的 Ferrari 预解三次在 \(K\) 上具有根 \(z\)。
四次多项式的经典判别指出：若预解三次在基域上有根，则伽罗瓦群稳定某一根配对，因而包含于配对稳定子 \(D_4\)。
进一步，取一个使特化伽罗瓦群为 \(D_4\) 的有理点 \(P\)，即可推出函数域伽罗瓦群必须达到 \(D_4\)。
最后由曲线覆盖的 Hilbert 不可约机制得到无穷多个保持该伽罗瓦群的有理点特化。证毕。
\end{proof}

\begin{proposition}[预解曲线上的二判别式锁定恒等式]\label{prop:xi-terminal-zm-dihedral-double-discriminant-lock}
\leanverified{paper\_xi\_terminal\_zm\_dihedral\_double\_discriminant\_lock}
在仿射曲线 \(\mathscr R(z,y)=0\) 上恒成立
\[
\boxed{\;
\bigl(z^{2}-4y(y+1)\bigr)\,\bigl(8y+5+4z\bigr)=(z+2)^{2}\; }.
\]
等价地，在多项式环 \(\ZZ[y,z]\) 中有恒等式
\[
\bigl(z^{2}-4y(y+1)\bigr)\,\bigl(8y+5+4z\bigr)-(z+2)^{2}=4\,\mathscr R(z,y).
\]
\end{proposition}

\begin{proof}
将左端展开并与 \((z+2)^2\) 相减，直接整理可得差为 \(4\mathscr R(z,y)\)。
因此在理想 \((\mathscr R)\) 上为零，等价于在曲线 \(\mathscr R=0\) 上恒成立。证毕。
\end{proof}

\begin{corollary}[Ferrari 分解的单二次域闭合与显式二次因子]\label{cor:xi-terminal-zm-dihedral-ferrari-single-quadratic-closure}
\leanverified{paper\_xi\_terminal\_zm\_dihedral\_ferrari\_single\_quadratic\_closure}
设 \(F\) 为特征不为 \(2\) 的域。
取 \(y,z\in F\) 满足 \(\mathscr R(z,y)=0\)，并在某个二次扩张中取
\[
s^2=z^2-4y(y+1).
\]
则在 \(F(s)[\lambda]\) 中存在显式因子分解
\[
\boxed{\;
\Pi(\lambda,y)=Q_{+}(\lambda)\,Q_{-}(\lambda)\;}
\]
其中
\[
Q_{\pm}(\lambda)=
\lambda^2-\frac12\left(1\pm\frac{z+2}{s}\right)\lambda+\frac12\bigl(z\pm s\bigr).
\]
并且由命题 \ref{prop:xi-terminal-zm-dihedral-double-discriminant-lock}，有
\[
\left(\frac{z+2}{s}\right)^2=8y+5+4z\in F,
\]
从而 Ferrari 过程中出现的两处平方根属于同一二次扩张 \(F(s)/F\)。
\end{corollary}

\begin{proof}[证明要点]
令 \(\lambda_1,\dots,\lambda_4\) 为 \(\Pi(\lambda,y)\) 的四根。
当 \(z\) 为预解三次的一根时，对应某一配对使得
\(
\lambda_1\lambda_2+\lambda_3\lambda_4=z
\)
且
\(
\lambda_1\lambda_2\lambda_3\lambda_4=y(y+1)
\)。
于是 \(p_{\pm}:=\lambda_1\lambda_2,\lambda_3\lambda_4\) 为二次方程 \(t^2-zt+y(y+1)=0\) 的两根，即 \(p_{\pm}=\frac12(z\pm s)\)。
另一方面，利用 \(\lambda_1+\cdots+\lambda_4=1\) 与二次对称函数关系可得两组和 \(\lambda_1+\lambda_2,\lambda_3+\lambda_4\) 的判别式为 \(8y+5+4z\)；
命题 \ref{prop:xi-terminal-zm-dihedral-double-discriminant-lock} 表明该判别式与 \(s^2\) 属于同一平方类，故同属 \(F(s)\)。
据此写出两二次因子并整理即得所示 \(Q_\pm\)。证毕。
\end{proof}

\begin{proposition}[\texorpdfstring{$D_4$}{D4} 特化时的二次子域生成元可计算性]\label{prop:xi-terminal-zm-dihedral-d4-quadratic-subfields}
\leanverified{paper\_xi\_terminal\_zm\_dihedral\_d4\_quadratic\_subfields}
设 \(y_0\in\QQ\) 使存在 \(z_0\in\QQ\) 满足 \(\mathscr R(z_0,y_0)=0\)，并且 \(\Pi(\lambda,y_0)\) 在 \(\QQ[\lambda]\) 上不可约且判别式非平方。
记其分裂域为 \(L/\QQ\)。
则 \(\mathrm{Gal}(L/\QQ)\cong D_4\)，并且 \(L\) 的二次子域可由两条平方根生成：
\[
K_1=\QQ\!\left(\sqrt{z_0^2-4y_0(y_0+1)}\right),
\qquad
K_2=\QQ\!\left(\sqrt{\mathrm{Disc}_\lambda\Pi(\lambda,y_0)}\right),
\]
第三个二次子域为 \(K_1K_2\) 的另一二次子域。
\end{proposition}

\begin{proof}[证明要点]
由 \(z_0\) 为预解根得 \(\mathrm{Gal}(L/\QQ)\subseteq D_4\)。
不可约性给出传递性，判别式非平方排除 \(V_4\)。
另一方面，推论 \ref{cor:xi-terminal-zm-dihedral-ferrari-single-quadratic-closure} 给出 \(K_1\subset L\)，而判别式二次分辨子域 \(K_2\subset L\) 恒成立。
若 \(\mathrm{Gal}\cong C_4\)，则分裂域仅含唯一二次子域，与 \(K_1\neq K_2\) 矛盾；故只能为 \(D_4\)。
二次子域格结构由 \(D_4\) 的指数 \(2\) 子群分类推出。证毕。
\end{proof}

\begin{conjecture}[沿无穷阶轨道的无例外二面体性]\label{conj:xi-terminal-zm-dihedral-d4-no-exception-orbit}
取结论 \ref{con:xi-terminal-zm-resolvent-cubic-mw-point} 的无穷阶点 \(P=(23,4)\in E_{\mathscr R}(\QQ)\)，并记 \(y_n:=y(nP)\in\QQ\)。
猜想对所有 \(n\ge 2\)，四次多项式 \(\Pi(\lambda,y_n)\) 在 \(\QQ\) 上不可约且
\[
\mathrm{Gal}\bigl(\Pi(\lambda,y_n)/\QQ\bigr)\cong D_4.
\]
\end{conjecture}
